\documentclass[11pt,a4paper]{article}
\usepackage[utf8]{inputenc}
\usepackage[T1]{fontenc}
\usepackage[italian]{babel}
\usepackage{amsmath,amssymb}
\usepackage[a4paper,margin=15mm]{geometry}
\usepackage{tikz}
\pagestyle{empty}
\setlength{\parindent}{0pt}
\setlength{\abovedisplayskip}{3pt}
\setlength{\belowdisplayskip}{3pt}
\newcommand{\taglio}{\par\noindent\tikz\draw[dashed,line width=.4pt] (0,0)--(\linewidth,0);\par}
\newcommand{\striscia}[5]{%
\taglio
\begin{minipage}[t][34mm][t]{\linewidth}
\vspace{2mm}\textbf{Scheda #1 -- #2}\hfill\textsf{#3}\par
\vspace{1mm}#4\par\vfill
{\small\textbf{Controllo precedente.} #5}\par\vspace{2mm}
\end{minipage}\par}
\newcommand{\inizio}[1]{%
\textbf{Laboratorio di equazioni -- Scheda #1}\par
{\small Lavora in coppia sul quaderno. Risolvi in $\mathbb R$; per A classifica l'equazione.
Nei quesiti teorici indica V/F e motiva. Ricevi una striscia alla volta.}\par\vspace{2mm}}
% Chiave docente: 1 indeterminata; 2 e 3 scomposizioni successive;
% 4 raccoglimento dell'espressione ripetuta e differenze di quadrati.
% I tagliandi finali non contengono esercizi e completano il controllo del n. 6.
\begin{document}
\inizio{1}
\striscia{1}{1}{A}{
\[\frac{3x-2}{\frac12+\frac14}-(x+1)=\frac{2x-1}{1-\frac13}-\frac{13}{6}\]
}{Inizio del percorso: nessun esercizio precedente.}
\striscia{1}{2}{B}{
\[x^4-10x^2+9=0\]
}{N. 1: $S=\mathbb{R}$; indeterminata.}
\striscia{1}{3}{C}{
\[x^5-2x^4-8x^3+16x^2+16x-32=0\]
}{N. 2: $S=\{-3;\,-1;\,1;\,3\}$.}
\striscia{1}{4}{D}{
\[\bigl[(x-1)^2-4\bigr]^2-12\bigl[(x-1)^2-4\bigr]=0\]
}{N. 3: $S=\{-2;\,2\}$.}
\striscia{1}{5}{Teoria -- equivalenza}{
\textbf{Vero o falso?} Aggiungendo $\frac1x$ a entrambi i membri di
$2x=0$ si ottiene un'equazione equivalente nell'insieme $\mathbb R$.
Motiva confrontando gli insiemi delle soluzioni.
}{N. 4: $S=\{-3;\,-1;\,3;\,5\}$.}
\striscia{1}{6}{Teoria -- classificazione}{
\textbf{Vero o falso?} L'equazione $0\cdot x=0$ ammette un'unica soluzione,
$x=0$. Motiva e indica la classificazione e l'insieme delle soluzioni.
}{N. 5: falso. Prima $S=\{0\}$; dopo $S=\{\}$, poich\'e $x=0$ \`e escluso.}
\taglio
\begin{minipage}[t]{\linewidth}
\vspace{2mm}{\small\textbf{Scheda 1 -- tagliando finale (consegnare dopo il n. 6).}\par
N. 6: falso. Ogni reale soddisfa l'equazione: indeterminata, $S=\mathbb{R}$.}\par\vspace{2mm}
\end{minipage}\par\taglio
\newpage
\inizio{2}
\striscia{2}{1}{A}{
\[\frac{3x-1}{\frac12+\frac14}-(x+2)=\frac{2x+1}{1-\frac13}-\frac{29}{6}\]
}{Inizio del percorso: nessun esercizio precedente.}
\striscia{2}{2}{B}{
\[x^4-13x^2+36=0\]
}{N. 1: $S=\mathbb{R}$; indeterminata.}
\striscia{2}{3}{C}{
\[x^5-3x^4-8x^3+24x^2+16x-48=0\]
}{N. 2: $S=\{-3;\,-2;\,2;\,3\}$.}
\striscia{2}{4}{D}{
\[\bigl[(x-2)^2-4\bigr]^2-12\bigl[(x-2)^2-4\bigr]=0\]
}{N. 3: $S=\{-2;\,2;\,3\}$.}
\striscia{2}{5}{Teoria -- equivalenza}{
\textbf{Vero o falso?} Sottraendo $\frac1{x-1}$ a entrambi i membri di
$3x=3$ si ottiene un'equazione equivalente nell'insieme $\mathbb R$.
Motiva confrontando gli insiemi delle soluzioni.
}{N. 4: $S=\{-2;\,0;\,4;\,6\}$.}
\striscia{2}{6}{Teoria -- classificazione}{
\textbf{Vero o falso?} L'equazione $0\cdot x=0$ non ammette soluzioni reali.
Motiva e indica la classificazione e l'insieme delle soluzioni.
}{N. 5: falso. Prima $S=\{1\}$; dopo $S=\{\}$, poich\'e $x=1$ \`e escluso.}
\taglio
\begin{minipage}[t]{\linewidth}
\vspace{2mm}{\small\textbf{Scheda 2 -- tagliando finale (consegnare dopo il n. 6).}\par
N. 6: falso. Ogni reale soddisfa l'equazione: indeterminata, $S=\mathbb{R}$.}\par\vspace{2mm}
\end{minipage}\par\taglio
\end{document}
